\begin{figure}[H]
\centering
\includegraphics[width=0.75\linewidth]{ehime/F8_diurnal.pdf}
\caption[日変化]{夜（03 JST）と昼（15 JST）の平均風速の差を、季節ごとに高度に対して描いたもの
（\SI{300}{\hecto\pascal} 以下の層のみ、領域平均）。横軸の目盛りに注意してほしい。
差は 3 領域・4 季節のどれでも \SI{1}{\metre\per\second} 未満で、
日々の変動（数十 \si{\metre\per\second}）に比べて小さい。
これが\textbf{時刻を区別せず平均してよい根拠}である。
ビンを時刻で 4 分割すると標本数も 4 分の 1 になるので、
必要のないところで分割しない理由がここにある。}
\label{fig:f8-ehime}
\end{figure}
